\clearpage
\TFMChapter{Metodología, Requisitos y Diseño de la Solución}

\TFMSection{Introducción}
\TFMSection{Plan de Trabajo y Metodología}
\TFMSection{Análisis de problema y caso de uso}
Por ejemplo:
\begin{tfmbullets}
\item Carga de documentación.
\item Procesamiento automático.
\item Inicio de una conversación.
\item Recopilación de requisitos del viaje.
\item Consulta de información turística.
\item Consulta de precios.
\item Elaboración del itinerario.
\item Generación de la propuesta comercial.
\end{tfmbullets}

\TFMSection{Requisitos Funcionales}
\TFMTOCPageBreak
\TFMSection{Requisitos No Funcionales}

\clearpage
\TFMSection{Arquitectura de la Aplicación}
Diagrama del flujo/funcionamiento. Presentar un diagrama completo con:
\begin{tfmbullets}
\item Frontend.
\item API backend.
\item Sistema multiagente.
\item Modelos de lenguaje.
\item RAG.
\item OpenSearch Serverless.
\item DynamoDB.
\item S3.
\item Lambdas.
\item Step Functions.
\item Servicios externos.
\end{tfmbullets}

\TFMSection{Diseño del sistema multiagente}
Explicar cada agente:
\begin{tfmbullets}
\item Agente orquestador.
\item Agente de información.
\item Agente de precios.
\item Agente de itinerario, si existe como agente independiente.
\item Agente generador de la propuesta comercial.
\end{tfmbullets}
También las transiciones y el estado compartido.

\TFMSection{Diseño del pipeline documental}
Desde que se carga un PDF hasta que queda disponible para las búsquedas:
\begin{enumerate}
\item Carga en S3.
\item Conversión de PDF a imágenes.
\item Extracción a Markdown.
\item Normalización.
\item Extracción de metadatos.
\item División en fragmentos.
\item Generación de embeddings o indexación.
\item Almacenamiento en OpenSearch.
\item Consulta posterior desde el RAG.
\end{enumerate}

Diseño de datos

\TFMSection{Diseño de datos}
DynamoDB, estructura de documentos en S3, índices de OpenSearch, metadatos, conversaciones y estado de las propuestas.

\TFMSection{Selección Tecnológica}
Justificar AWS, OpenAI, LangGraph, OpenSearch Serverless, Lambda, Step Functions, DynamoDB, frontend y backend utilizados.
Yo movería Selección tecnológica después de describir la arquitectura, porque así primero explicáis qué necesita el sistema y luego justificáis las tecnologías elegidas.
